\documentclass[conference]{IEEEtran}
\IEEEoverridecommandlockouts

% ─── Packages ────────────────────────────────────────────────────────────────
\usepackage{cite}
\usepackage{amsmath,amssymb,amsfonts}
\usepackage{algorithmic}
\usepackage{algorithm}
\usepackage{graphicx}
\usepackage{textcomp}
\usepackage{xcolor}
\usepackage{booktabs}
\usepackage{multirow}
\usepackage{url}
\usepackage{hyperref}
\usepackage{microtype}
\usepackage{enumitem}
\usepackage[T1]{fontenc}
\usepackage[utf8]{inputenc}
\usepackage{listings}
\usepackage{array}
\usepackage{balance}
\usepackage{titlesec}
% ─── Listings style ──────────────────────────────────────────────────────────
\lstset{
  basicstyle=\ttfamily\scriptsize,
  breaklines=true,
  frame=single,
  numbers=left,
  numberstyle=\tiny,
  keywordstyle=\color{blue},
  commentstyle=\color{gray},
  stringstyle=\color{red!60!black}, 
}

% ─── Hyperref config ─────────────────────────────────────────────────────────
\hypersetup{
  colorlinks=true,
  linkcolor=blue!70!black,
  citecolor=green!50!black,
  urlcolor=blue!70!black,
}

\def\BibTeX{{\rm B\kern-.05em{\sc i\kern-.025em b}\kern-.08em
    T\kern-.1667em\lower.7ex\hbox{E}\kern-.125emX}}

% ─────────────────────────────────────────────────────────────────────────────

\begin{document}

\title{Engineering a Hallucination-Resistant, Multilingual RAG Architecture
       for Indian Parliamentary Data}

\author{%
  \IEEEauthorblockN{Vivek Rekha Ashoka}
  \IEEEauthorblockA{%
    \textit{Department of Computer Science \& Engineering}\\
    \textit{PES University}\\
    \textit{Bangalore, India}\\
    pes1ug23cs704@pesu.pes.edu}
\and
  \IEEEauthorblockN{Vrishant Bhalla}
  \IEEEauthorblockA{%
     \textit{Department of Computer Science \& Engineering}\\
    \textit{PES University}\\
    \textit{Bangalore, India}\\
    pes1ug23cs706@pesu.pes.edu}
}

\maketitle

% ─────────────────────────────────────────────────────────────────────────────
\begin{abstract}
When implementing a RAG pipeline in a legal/institutional domain where it is very important to be factually accurate, some common failures that can occur are:
(i)~Large Language Model (LLM) hallucinations that generate fictitious statutory citations;
(ii)~the loss of precision in numbers resulting from sparse keyword models failing to generalise synonyms of legal phrasing; and
(iii)~degraded retrieval and reranking quality when the corpus contains documents in multiple languages and scripts, causing tokenisation-induced truncation in cross-encoders optimised for a single language.
We evaluate the effectiveness of our system using statutory fact queries from the Lok Sabha Q\&A data source for the 16th to 18th Lok Sabha.
We have developed a two-stage hybrid retrieval system, under hardware constraints, that aims to rectify each of these failure modes using a principled engineering pipeline that includes: a multilingual semantic chunking strategy designed with regard to the structure of the statutory data; fusing of \textit{BAAI/bge-m3} dense vectors and BM25 sparse arrays, stored and indexed on Qdrant, with the use of Reciprocal Rank Fusion (RRF) for initial retrieval followed by the use of a \textit{BAAI/bge-reranker-base} cross-encoder to implement re-ranking of the retrieved document collections; and the use of strictly constrained Cloud LLMs to generate the desired output by using retrieved documents, which provides a desired output according to its prompt guidelines.
\end{abstract}

\begin{IEEEkeywords}
Retrieval-Augmented Generation, Hybrid Retrieval, Reciprocal Rank Fusion,
Cross-Encoder Reranking, Multilingual NLP, Hallucination Mitigation, Qdrant,
BM25, Parliamentary AI, Indian Legal NLP, Hardware-Constrained Deployment
\end{IEEEkeywords}

% =============================================================================
\section{Introduction}
\label{sec:intro}

\subsection{Problem Context and Motivation}

The Lok Sabha, the lower house of the Indian Parliament, generates an
extraordinary volume of legislative artefacts spanning several decades:
plenary debate transcripts, Starred and Unstarred question-answer sessions,
Bills, Acts, committee reports, budget speeches, and ministerial statements.
These documents are heterogeneous in structure, language, and numerical
density; collectively they constitute a corpus of $\approx$ 200,000 semantically
chunked text fragments once pre-processed for machine ingestion.

The emergence of powerful LLMs has generated significant interest in deploying
conversational AI interfaces atop such corpora~\cite{lewis2020rag}.
In legal and parliamentary contexts, however, the cost of an incorrect
response is qualitatively different from a factual error in a general-purpose
chatbot.
A fabricated statutory penalty, a misquoted Bill number, or an invented
committee recommendation can mislead legislators, journalists, and the public
in consequential ways~\cite{rag_framework2025,kemenkeu2024}.
Base LLMs, trained to produce fluent text rather than verified citations,
routinely hallucinate with high confidence when queried on low-frequency
factual content~\cite{lewis2020rag}, making unguarded deployment over
parliamentary data ethically untenable.

Beyond hallucination, the linguistic profile of the corpus compounds
retrieval difficulty.
Hindi legislative documents co-exist with English documents, and standard
English-optimised cross-encoders suffer from tokenisation-induced degradation
on Devanagari and Dravidian scripts~\cite{multilingualrag_2025,mrag_settings2024}.
No single retrieval modality---dense semantic vectors or sparse keyword
arrays---handles the full breadth of statutory precision queries on its own.

\subsection{Objective and Contributions}

The primary objective of this research is to architect, deploy, and evaluate
a \textbf{hallucination-resistant, hardware-optimised Two-Stage Hybrid RAG
System} for the Lok Sabha parliamentary corpus.
Concretely, the contributions are:

\begin{enumerate}[leftmargin=*, label=(\roman*)]
  \item An end-to-end data pipeline for constructing a structured,
    machine-readable parliamentary corpus from scratch---including PDF
    scraping, text extraction, boilerplate removal, failure cataloguing,
    and structural chunking---over the 16th--18th Lok Sabha archives.
  \item A structural regex chunker with a 400-word ceiling and
    parent-context propagation, addressing the ``Blind Split'' problem for
    both Q\&A and Debate/Bill document types.
  \item A multilingual ingestion and retrieval strategy that combines a
    language-aware embedding encoder (BAAI/bge-m3, trained on 100+ languages)
    with a multilingual cross-encoder reranker (BAAI/bge-reranker-v2-m3),
    mitigating tokenisation-induced truncation on Devanagari-script documents.
  \item An architectural configuration recipe for deploying a
    $\approx$200,000-document Qdrant vector database on consumer hardware
    (16\,GB RAM) while avoiding OOM crashes, using on-disk vector and HNSW
    storage.
\end{enumerate}

% =============================================================================
\section{Background and Literature Review}
\label{sec:background}

\subsection{Key Technical Concepts}

\subsubsection{Dense Embeddings vs.\ Sparse Keyword Arrays}

\emph{Dense embeddings} encode text as continuous vectors
$\mathbf{v} \in \mathbb{R}^d$ via transformer bi-encoders; cosine similarity
enables semantic retrieval even when surface vocabulary differs~\cite{lewis2020rag}.
This work uses BAAI/bge-m3 ($d=1024$, 100+ languages)~\cite{mrag_settings2024}.

\emph{Sparse keyword arrays} (BM25) weight terms via inverted-index vectors.
The relevance score is:

\begin{equation}
  \text{BM25}(D,Q) = \sum_{i=1}^{n} \text{IDF}(q_i) \cdot
  \frac{f(q_i,D)\,(k_1\!+\!1)}{f(q_i,D)+k_1\!\left(1-b+b\frac{|D|}{\overline{dl}}\right)}
  \label{eq:bm25}
\end{equation}

\noindent where $f(q_i, D)$ is term frequency, $|D|$ document length,
$\overline{dl}$ average document length, $k_1{\approx}1.2$--$2.0$, and
$b=0.75$~\cite{rag_framework2025}.
BM25 excels at exact-match recall for statutory section numbers and Bill
identifiers, but its unbounded, corpus-dependent scores make direct fusion
with bounded cosine similarities problematic.

\subsubsection{Evolution from Naive to Advanced RAG}

\emph{Naive RAG}~\cite{lewis2020rag} encodes a query with a bi-encoder and
performs approximate nearest-neighbour (ANN) search.
While computationally efficient, bi-encoders embed query and document
independently, so relevance is assessed in an abstract latent space without
direct attention to query--document token interactions.

\emph{Advanced RAG} introduces a second-stage cross-encoder reranker that
receives the concatenated query--document pair as a single input~\cite{mrag_settings2024},
enabling full self-attention across the pair and yielding substantially higher
precision at the cost of $O(K)$ additional inference calls, where $K$ is the
candidate set size from Stage~1.

\subsubsection{Lok Sabha Data Structure}

The corpus processed in this work spans the 16th, 17th, and 18th Lok Sabhas
and comprises Starred Questions, Unstarred Questions, Bills, and Debates in PDF format.
Documents exhibit extreme structural heterogeneity: short single-paragraph
Q\&A pairs of fewer than 100 words co-exist with multi-chapter Acts spanning
hundreds of pages.
English-language structural markers such as \texttt{CHAPTER}, \texttt{SECTION},
\texttt{ANNEXURE}, and \texttt{APPENDIX} are present alongside
Hindi-language parliamentary content, making bilingual chunking a prerequisite
for faithful document segmentation.

\subsection{Literature Review}

\subsubsection{Government and Legal RAG Applications}

KemenkeuGPT~\cite{kemenkeu2024} provides the closest parallel to this work:
an iterative RAG system built over Indonesia's Ministry of Finance regulations
and financial data spanning 2003--2023 using the LangChain framework.
Using the RAGAS evaluation framework, KemenkeuGPT achieved 44\% correctness
with 73\% faithfulness on government regulatory queries.
Critically, the authors identify that model accuracy improved from 35\% to
61\% through iterative prompt engineering and fine-tuning, demonstrating the
importance of generation-layer constraints in government document RAG.
Unlike our system, however, KemenkeuGPT operates mono-lingually over
Indonesian text and does not address the cross-lingual retrieval challenge.

Graph-based RAG extensions~\cite{graphrag2025,tkgrag2025} and Knowledge
Graph-guided approaches~\cite{kgrag2025} have been proposed to capture
entity relationships across document boundaries---a capability directly
relevant to mapping legislative amendment chains across Lok Sabha sessions.
While these approaches offer richer inter-document reasoning, their
construction complexity and storage overhead are incompatible with the
consumer-hardware constraints of the present deployment.

\subsubsection{Multilingual RAG Systems}

Chirkova and Nikoulina~\cite{mrag_settings2024} establish the first
comprehensive baseline for multilingual RAG (mRAG) across 13 languages
using the BERGEN benchmarking library.
Their key findings directly inform this work:
(i)~Off-the-shelf multilingual retrievers including BAAI/bge-m3 perform
reasonably well in cross-lingual settings; (ii)~achieving high performance
across non-English queries requires a strong multilingually pre-trained LLM
coupled with advanced prompting strategies specifying the desired output
language; and (iii)~code-switching in non-Latin-alphabet languages (including
Devanagari) remains a principal failure mode of current multilingual generation
systems.
Their work notes that ``for each question in the evaluation data, we retrieve
50 relevant passages and pass them to the reranker to select top-5 relevant
ones''---a pipeline architecture that directly mirrors the prefetch=50,
rerank-to-top-$k$ design adopted in this work.

Ranaldi et al.~\cite{multilingualrag_2025} propose CrossRAG, which translates
retrieved documents to a common language (English) before generation.
Their findings show that CrossRAG ``significantly enhances performance on
knowledge-intensive tasks, benefiting both high-resource and low-resource
languages.''
The present work addresses the same cross-lingual challenge not through
post-retrieval translation (which incurs per-query API latency) but through
pre-ingestion language tagging and the selection of a retrieval encoder
(bge-m3) that natively handles multilingual similarity in a single embedding
space.

% =============================================================================
\section{System Architecture, Methodology, and Implementation}
\label{sec:arch}

\subsection{Scope and Operational Constraints}

The system operates over a \emph{static} pre-chunked dataset of $\approx$
\textbf{200,000 JSON objects} covering the Q\&A sessions and Debate/Bill
documents of the 16th, 17th, and 18th Lok Sabhas.
No live web scraping or real-time parliamentary feed integration is
implemented.

Consumer-hardware constraints are fundamental architecture drivers.
The target machine provides a 16\,GB unified RAM envelope with no discrete GPU,
obligating all embedding and reranking inference to run via CPU backends
(PyTorch CPU with ONNX where available).
Cloud API quotas on the Google Gemini free-tier (approximately 15 requests per
minute) constrain the generation layer and motivate the offline/online
processing split: all embedding, chunking, and language tagging are performed
offline and stored persistently; only final-stage generation is performed
online at query time.

\subsection{End-to-End Pipeline Overview}

The complete system consists of four offline pipeline stages and two online
query-time stages:

\begin{enumerate}[leftmargin=*, label=\textbf{Stage \Roman*.}]
  \item \textbf{(Offline) Document Extraction \& Language Tagging}:
    PDF extraction via PyMuPDF; boilerplate removal; \texttt{langdetect}
    language annotation per document.
  \item \textbf{(Offline) Structure-Aware Chunking \& Context Injection}:
    Regex-based structural splitting and parent-context prepending; output
    to JSON corpus files.
  \item \textbf{(Offline) ID Namespacing \& Deduplication}:
    Prefix-based corpus merging (\texttt{QA\_}, \texttt{DB\_}) and
    sub-chunk collision resolution.
  \item \textbf{(Offline) Vector Ingestion}:
    SQLite text vault population and Qdrant dual-vector upload.
  \item \textbf{(Offline) Stage V -- Hybrid RRF Retrieval}:
    Dual-path prefetch (dense + sparse) fused via RRF; top-50 candidates.
  \item \textbf{(Online) Stage VI -- Cross-Encoder Reranking \& Generation}:
    Reranker-scored top-$k$ reduction; Gemini-2.5-Flash constrained
    generation.
\end{enumerate}

\begin{figure}[t]
    \centering
    \includegraphics[width=\columnwidth]{image.png}
    \caption{End-to-end layered architecture of the Two-Stage Hybrid RAG pipeline.}
    \label{fig:architecture}
\end{figure}

\subsection{Experimental Setup and Hyperparameter Configuration}

\begin{table}[t]
  \caption{Core Hyperparameters of the Two-Stage Hybrid RAG System}
  \label{tab:hyperparams}
  \centering
  \footnotesize
  \begin{tabular}{@{}p{2.2cm}p{2.4cm}p{2.4cm}@{}}
    \toprule
    \textbf{Parameter} & \textbf{Value} & \textbf{Rationale} \\
    \midrule
    Chunk word limit   & 400 words       & 512-token budget \\
    Dense encoder      & BAAI/bge-m3     & 100+ lang.; 1024D \\
    Sparse encoder     & Qdrant/bm25     & Exact keyword match \\
    Reranker (v1)      & bge-reranker-base & Fast; English CE \\
    Reranker (v2)      & bge-reranker-v2-m3 & Multilingual CE \\
    Prefetch $K$       & 50 docs         & Both retrieval paths \\
    Final $k$          & 10 docs         & LLM context budget \\
    RRF constant       & Default (Qdrant)& Server-side fusion \\
    Vector dim.        & 1024            & bge-m3 output \\
    Qdrant storage     & On-disk SSD     & OOM prevention \\
    CPU threads        & 4               & Context-switch guard \\
    LLM backbone       & Gemini Flash    & Low latency \\
    CE batch size      & 4 pairs         & CPU memory budget \\
    \bottomrule
  \end{tabular}
\end{table}

Table~\ref{tab:hyperparams} summarises the key hyperparameters.
The software stack comprises Python~3.x, Qdrant~(Docker), HuggingFace
\texttt{sentence-transformers}, \texttt{fastembed}~(for BM25 sparse vectors
via the \texttt{Qdrant/bm25} model), PyMuPDF~(\texttt{fitz}), \texttt{langdetect},
SQLite3, and the Google GenAI SDK.

\subsection{Stage 0: Document Extraction, Cleaning, and Language Tagging}
\label{sec:extraction}

\subsubsection{PDF Extraction via PyMuPDF}

Parliamentary Q\&A documents in the 16th--18th Lok Sabha archives are
distributed as individual PDFs.
The \texttt{QAExtractor.py} script uses PyMuPDF~(\texttt{fitz}) to iterate
over each page and extract digital text via \texttt{page.get\_text("text")}.
Documents with fewer than 50 characters of extracted text are discarded as
scanned images or corrupt files; the \texttt{FailureAnalyzer.py} script
catalogues such failures into three categories: \emph{Missing File},
\emph{Corrupted/HTML}, and \emph{Scanned Image (No Text)}.

\subsubsection{Boilerplate Removal}

A multi-pattern regex stripping pass removes parliamentary boilerplate that
would otherwise inflate term frequencies and bias BM25 scoring.

\subsubsection{Language Detection and Metadata Tagging}

After cleaning, each document is passed to \texttt{langdetect.detect()},
which returns an ISO~639-1 language code stored under the \texttt{language}
key in the per-document metadata JSON:


\noindent The \texttt{DetectorFactory.seed = 0} call enforces deterministic
language detection results across runs.
This language tag is stored in the Qdrant payload alongside \texttt{document\_id},
enabling future language-stratified retrieval or filtered queries without
re-running detection at query time.

\subsection{Stage 1: Structure-Aware Chunking and Context Propagation}
\label{sec:chunking}

\subsubsection{The Blind Split Problem}

Naive chunking by fixed character or word count cleaves document structure
arbitrarily, producing orphaned child chunks that lack the document heading
or section context required for accurate embedding.
A chunk containing only a table of financial penalties---without the enclosing
section title---may not retrieve for a query referencing that section by number,
because the section number resides in an adjacent sibling chunk.

\subsubsection{Structural Regex Chunker}

The chunker (\texttt{QAChunker.py}, \texttt{DebBillChunker.py}) enforces a
hard ceiling of \texttt{WORD\_LIMIT = 400} words per chunk, which was
selected to remain safely below the 512-token limit of the downstream
cross-encoder models (approximately 380 words of English text at average
subword tokenisation rates).
Documents already below 400 words pass through unchanged.

Long documents are split by a positive-lookahead structural regex that
preserves the trigger header word at the start of each child chunk:


If the regex finds no structural boundary (a continuous wall of text), the
fallback splits on two or more consecutive newlines (\texttt{r'\textbackslash
n\{2,\}'}), simulating a paragraph boundary.

\subsubsection{Parent-Context Injection}

A child chunk at index~$> 0$ within a split document receives an explicit
context header prepended to its text before embedding:

This single concatenation restores the document-level discriminative signal
that a bi-encoder would otherwise lack.
Child chunk IDs are suffixed with \texttt{\_part\{index\}} to maintain
uniqueness within the SQLite primary key constraint.

\subsection{Stage 2: ID Namespacing and Deduplication}

The corpus is split across two sub-teams producing two JSON files:
\texttt{patched\_corpus\_A.json} (Q\&A documents, prefix \texttt{QA\_}) and
\texttt{final\_corpus\_B.json} (Debate/Bill documents, prefix \texttt{DB\_}).
\texttt{CorpusPatcher.py} prepends the appropriate prefix to all document
IDs, and \texttt{DebBillDupeFixer.py} resolves any remaining ID collisions
by appending an incrementing \texttt{\_subchunk\{N\}} suffix.
The resulting IDs are globally unique across the merged $\approx$200,000-document
corpus, ensuring that both the SQLite primary key constraint and the Qdrant
UUID5-based point ID generation produce collision-free entries.

\subsection{Stage 3: Vector Ingestion -- The Dual-Store Architecture}
\label{sec:ingestion}

\subsubsection{The SQLite Text Vault}

Full raw text is stored in a local SQLite database (\texttt{loksabha\_text\_store.db})
with the schema:

\begin{lstlisting}[language=SQL,
  caption={SQLite documents table schema.}]
CREATE TABLE IF NOT EXISTS documents (
    chunk_id  TEXT PRIMARY KEY,
    doc_type  TEXT,
    title     TEXT,
    raw_text  TEXT
);
\end{lstlisting}

The \texttt{INSERT OR IGNORE} idiom allows idempotent multi-corpus merging:
documents already present (matched on \texttt{chunk\_id}) are silently
skipped, enabling the two sub-corpora to be inserted in either order without
creating duplicates.
Qdrant point payloads store only lightweight metadata (\texttt{document\_id},
\texttt{doc\_type}, \texttt{title}, \texttt{language}, etc.);
the full text is retrieved from SQLite at query time, decoupling the
memory-mapped vector index from bulky text blobs.

\subsubsection{Qdrant Collection Configuration: The Low-RAM Straitjacket}
\label{sec:qdrant-config}

Qdrant by default pins vectors and HNSW graph edges in RAM.
At $d=1024$ with \texttt{float32} precision, the raw vector footprint alone is:

\begin{equation}
  M_\text{vec} = 200{,}000 \times 1024 \times 4\,\text{B}
  \approx 986\,\text{MB}
  \label{eq:ramvec}
\end{equation}

With HNSW graph overhead, payload dictionaries, and concurrent segment-merge
RAM spikes, the in-memory footprint was observed to exceed available RAM
during pilot ingestion, causing kernel OOM kills.

The two \texttt{on\_disk=True} flags instruct Qdrant to memory-map all
vector data and HNSW graph edges from the SSD rather than pinning them in
RAM.
Page-fault overhead is negligible at query time because only the $K=50$
prefetched vector records need to be loaded per query.
Ingestion is performed in batches of 1,000 points (\texttt{BATCH\_SIZE=1000}),
and each point receives a deterministic UUID5 derived from its corpus chunk ID
to ensure idempotent re-uploads.

The hybrid Qdrant collection (\texttt{loksabha\_rag\_hybrid\_bm25\_bge})
used at query time stores both a \texttt{dense\_bge} named dense vector
namespace and a \texttt{sparse\_keyword} named sparse vector namespace,
enabling the dual-path prefetch query described in Section~\ref{sec:stage1}.

\subsection{Stage 4 (Query-Time): Hybrid RRF Retrieval}
\label{sec:stage1}

\subsubsection{Dense Retrieval -- BAAI/bge-m3}

The primary retrieval encoder is BAAI/bge-m3~\cite{mrag_settings2024},
a multilingual bi-encoder producing 1024-dimensional vectors trained on
100+ languages.
Query encoding is performed via \texttt{SentenceTransformer.encode()},
which delegates to a CPU ONNX backend when no CUDA device is detected.

The \texttt{torch.set\_num\_threads(4)} call is critical on consumer laptops:
without this constraint, PyTorch spawns threads equal to the number of logical
CPU cores, which causes excessive context-switching overhead on a shared OS
workload.

\subsubsection{Sparse Retrieval -- BM25 via FastEmbed}

The secondary retrieval path uses the \texttt{Qdrant/bm25} model loaded
via the \texttt{fastembed} library's \texttt{SparseTextEmbedding} class.
BM25 sparse vectors are produced at both ingestion time (for the document
index) and query time (for the query) using the same tokeniser, ensuring
vocabulary alignment.
The resulting sparse vector is represented as a pair of parallel arrays
(\texttt{indices}, \texttt{values}) and wrapped in a Qdrant
\texttt{SparseVector} object.

The lexical path is essential for the Lok Sabha corpus because statutory
section numbers (e.g., ``Section~82''), monetary fine amounts, and Bill
identifiers are high-specificity, low-frequency terms for which dense
similarity may retrieve thematically related but numerically incorrect passages.

\subsubsection{Score Fusion: Why RRF over Direct Score Averaging}

Two fusion strategies were considered at design time: (a)~\emph{direct linear
combination} of normalised scores, and (b)~\emph{Reciprocal Rank Fusion}
(RRF).

\textbf{Direct score combination} requires normalising the raw BM25 scores,
which are unbounded above.
Because BM25 scores are corpus-dependent and query-dependent, a fixed
normalisation range cannot be determined without running the full corpus
query, making linear combination impractical for runtime deployment.
An outlier document that matches a repetitive parliamentary formulae many times
produces a BM25 score orders of magnitude above average, collapsing the
effective range of other documents after normalisation.

\textbf{RRF} avoids score magnitude entirely by operating on \emph{ranks}.
For retrieval lists $L_1$ (dense) and $L_2$ (sparse), each returning up to
$K=50$ candidates, the RRF score for document $d$ is:

\begin{equation}
  \text{RRF}(d) = \sum_{m \in \{L_1, L_2\}} \frac{1}{k + r_m(d)}
  \label{eq:rrf}
\end{equation}

\noindent where $r_m(d)$ is the 1-indexed rank of $d$ in list $L_m$ (missing
documents contribute zero), and $k$ is a smoothing constant.
A document appearing at rank~3 in dense retrieval and rank~8 in BM25 receives
$\frac{1}{k+3}+\frac{1}{k+8}$, reliably outranking a document present in only
one list regardless of the magnitude of its raw scores.

Qdrant implements RRF server-side as a named fusion strategy
(\texttt{models.Fusion.RRF}), invoked through the \texttt{FusionQuery}
prefetch mechanism:

The \texttt{limit=50} on the outer query returns the top-50 documents from
the RRF-merged ranking, providing the Stage-2 reranker with its full
candidate pool.

\subsection{Stage 5 (Query-Time): SQLite Text Retrieval and Cross-Encoder Reranking}
\label{sec:stage2}

\subsubsection{Bulk SQLite Fetch -- Eliminating the N+1 Problem}

Version~1 (\texttt{SearchGenerateReRank.py}) retrieves document text by
executing 50 individual \texttt{SELECT} queries---one per Qdrant hit.
This N+1 pattern introduces substantial I/O overhead.
Version~2 (\texttt{SearchGenerateReRank-v2.py}) replaces this with a single
bulk \texttt{IN} query.

The subsequent per-document lookup in \texttt{doc\_map} is $O(1)$, reducing
50 sequential disk reads to a single query against the SQLite B-tree index.

\subsubsection{Cross-Encoder Reranking}

\textbf{Version 1} loads \texttt{BAAI/bge-reranker-base}, a BERT-class
12-layer cross-encoder fine-tuned on MS-MARCO~\cite{mrag_settings2024}.
\textbf{Version 2} upgrades to \texttt{BAAI/bge-reranker-v2-m3}, the
multilingual variant of the same architecture that handles non-Latin scripts
more robustly at the cost of higher inference latency.

In both versions, the 50 candidate documents retrieved from Stage~4 are
formed into (query, document\_text) pairs and scored by the cross-encoder:

\noindent The \texttt{batch\_size=4} parameter limits the number of
(query, document) pairs processed simultaneously in a single forward pass,
preventing host RAM spikes from concurrent BERT activations on CPU.

The cross-encoder outputs raw logit scores reflecting the joint relevance of
each query--document pair.
Because both query and document tokens attend to each other symmetrically at
every encoder layer, the model captures fine-grained lexical and syntactic
interactions that are impossible in the independent bi-encoder formulation.
The $k$-highest scoring documents are sliced from the sorted list and
passed to the generation stage.

The English-optimised \texttt{bge-reranker-base}
operates within a fixed 512-token window.
Since Devanagari text consumes significantly more tokens per word than English
under subword tokenisers trained predominantly on Latin-script data, a Hindi
document of moderate word length can silently exceed 512 tokens, causing the
cross-encoder to truncate and score an incomplete passage.
Selecting BAAI/bge-m3 as the bi-encoder (multilingual, 100+ languages) for
Stage~1 means that semantically relevant Hindi chunks \emph{will} appear in
the top-50 candidate list.
Version~2 resolves this by replacing the reranker with
\texttt{bge-reranker-v2-m3}, which was trained on multilingual data and
handles Devanagari tokenisation more accurately, at the cost of increased
CPU inference time.

\subsection{Stage 6: Constrained Generation}
\label{sec:generation}

The top-$k$ reranked documents are serialised into a structured context block
with source attribution headers and submitted to
\texttt{gemini flash} via the Google GenAI SDK

\subsubsection{Source-Constrained Prompt Protocol}

The system prompt enforces three non-negotiable rules:

\begin{lstlisting}[language=Python,
  caption={Generation prompt with hallucination-prevention rules.}]
prompt = 
f"""You are a precision-focused legal and parliamentary assistant for the Government of India.
Your objective is to synthesize accurate, highly structured answers to the user's query using ONLY the provided Lok Sabha records.
CORE DIRECTIVES:
 1. ABSOLUTE GROUNDING: You must rely exclusively on the CONTEXT provided below. Do not use external knowledge, internet training data, or infer legal facts not explicitly stated.
 2. THE "PIVOT & INFORM" RULE: If the CONTEXT does not contain the exact information required to fully answer the query, you MUST begin your response by explicitly stating: "I cannot directly answer this specific question based on the provided parliamentary records." Immediately following this, you must pivot to provide a structured summary of the most closely related legal facts, debates, or definitions that *are* present in the retrieved CONTEXT.
 3. CITATION MANDATE: Every factual claim, statistic, or penalty in your answer (including the related information provided under Rule 2) MUST be immediately followed by an inline citation to its source document (e.g., [SOURCE: BILL - The Personal Data Protection Bill, 2019]). 
 4. CHRONOLOGY & CONFLICTS: If the CONTEXT contains multiple documents showing how a law or debate evolved over time (e.g., a 2019 draft vs. a 2023 amendment), clearly separate them and explain the evolution. 
 5. STRUCTURAL CLARITY: Output your response using strict Markdown. Use bullet points for lists, and use bold text for specific penalties, dates, and Section/Clause numbers. Maintain a formal, objective, and non-partisan tone.

CONTEXT:
{context_string}

USER QUESTION: 
{user_question}

FINAL ANSWER:"""
\end{lstlisting}

\noindent Rule~1 instructs the LLM to \emph{refuse} rather than fabricate
when the retrieved context is insufficient.
Rule~3 binds every factual claim to a named source, providing an audit trail.
Together these rules implement the ``closed-book'' generation constraint that
is essential for a parliamentary fact-retrieval system.

% =============================================================================
\section{Results and Evaluation}
\label{sec:results}

\subsection{Architecture Evolution and Iteration}

The system was developed through three measurable architectural iterations
summarised in Table~\ref{tab:evolution}.

\begin{table}[t]
  \caption{Architectural Iterations of the Lok Sabha RAG Pipeline}
  \label{tab:evolution}
  \centering
  \footnotesize
  \renewcommand{\arraystretch}{1.3}
  \begin{tabular}{@{}p{0.7cm}p{3.0cm}p{3.3cm}@{}}
    \toprule
    \textbf{Ver.} & \textbf{Key Configuration} & \textbf{Primary Change} \\
    \midrule
    v0 & Dense-only retrieval;
         \texttt{SearchTest.py}
       & Baseline ANN; no fusion, no reranker \\
    \addlinespace
    v1 & RRF + \texttt{bge-reranker-base};
         N+1 SQLite fetch
       & Two-stage hybrid pipeline; English cross-encoder \\
    \addlinespace
    v2 & RRF + \texttt{bge-reranker-v2-m3};
         bulk SQL; thread cap
       & Multilingual CE; N+1 fix; CPU optimisation \\
    \bottomrule
  \end{tabular}
  \renewcommand{\arraystretch}{1.0}
\end{table}

\subsection{Impact of Prefetch Depth on Recall}

A critical architectural decision is the Stage-1 prefetch depth $K$.
Setting $K$ too low causes split legal clauses to be missed: if a document's
section heading and its penalty sub-clause reside in adjacent chunks, both
must appear in the top-50 candidates for the cross-encoder to surface the
sub-clause.
The code comment in \texttt{SearchGenerateReRank.py} explicitly motivates this
choice: \textit{``We fetch 50 documents instead of 10 to give the reranker
more options''}, highlighting that the 50-document prefetch was a deliberate
engineering decision to recover split legal context.

\subsection{Cross-Encoder and Multilingual Trade-off}

The upgrade from \texttt{bge-reranker-base} (v1) to
\texttt{bge-reranker-v2-m3} (v2) directly deploys
a multilingual cross-encoder that is not constrained by English subword
tokenisation for Devanagari input.
The \texttt{batch\_size=4} setting in the v2 reranker constrains peak
memory consumption during the 50 forward passes to a level compatible with
the 16\,GB envelope.
The trade-off is increased per-query reranking latency relative to v1,
which is acceptable given the offline-vs-online processing split.

\subsection{Hallucination Mitigation via Prompt Constraints}

The three-rule prompt constraint follows a design principle validated in
government LLM deployments~\cite{kemenkeu2024}: explicitly grounding the
model in provided context and mandating source citations substantially
reduces fabrication.
KemenkeuGPT~\cite{kemenkeu2024} demonstrated that iterative prompt engineering
on a government RAG system improved correctness from 48\% to 64\%.
The present system adopts an analogous closed-book constraint: any claim not
directly supported by a named source chunk must be actively refused rather
than guessed.

\subsection{Qualitative Behaviour}

For queries involving statutory penalties or financial figures, the source
attribution format \texttt{[SOURCE: BILL - <title>]} allows users to
immediately verify the cited document.
When the retrieved context contains a query-relevant section whose
numeric value is absent or truncated (a common artefact of PDF
extraction), the LLM is expected to produce ``I cannot answer this based on
the provided parliamentary records'' rather than hallucinating a plausible
number---consistent with Rule~1 of the prompt protocol.

This behaviour is structurally analogous to the \emph{faithfulness} metric
measured in KemenkeuGPT~\cite{kemenkeu2024} using RAGAS (73\% faithfulness
achieved in that work), though a comparative RAGAS evaluation on the Lok Sabha
corpus is deferred to future work due to the absence of a ground-truth
annotation set.
Quantitative metrics including RAGAS faithfulness, answer correctness,
and context precision over a human-annotated evaluation set remain as
\textbf{[INSERT QUANTITATIVE EVALUATION RESULTS HERE]} pending the construction
of a labelled benchmark.

% =============================================================================
\section{Discussion and Limitations}
\label{sec:discussion}

\subsection{The API Dependency Bottleneck}

The generation layer depends on live access to the Google Gemini cloud API.
At free-tier rate limits, the system is constrained to approximately
15~generation requests per minute.
A Gemini service outage renders the generation stage unavailable; however,
the retrieval pipeline (Stages~4--5) remains fully functional offline, and
the top-$k$ reranked document set can be returned directly to a user
interface as a ranked evidence list even without LLM generation.
A production deployment would require either a rate-limited retry queue or
a fallback local generation model (e.g., a quantised Llama variant).

\subsection{Multilingual Reranker Latency}

The \texttt{bge-reranker-v2-m3} model is substantially larger than
\texttt{bge-reranker-base} and requires more CPU computation per forward
pass.
On a consumer laptop without a GPU, the latency of 50 cross-encoder forward
passes at \texttt{batch\_size=4} may become a bottleneck for interactive
query response times.
The \texttt{batch\_size=4} setting was tuned to balance memory safety against
throughput; increasing it risks RAM pressure, while decreasing it further
adds sequential I/O overhead.

\subsection{BM25 Stop-Word Inefficiencies on Non-English Text}

The \texttt{Qdrant/bm25} model implements an English stop-word list.
Hindi function words that are not in the English stop-word list receive
non-zero IDF weights and contribute spurious BM25 signal.
Prior to language-tag filtering, this can cause Hindi documents with
high-frequency Hindi function words to accumulate inflated BM25 scores
relative to structurally equivalent English documents.
The \texttt{language} metadata tag stored at ingestion time (from
\texttt{langdetect}) enables downstream filtering to mitigate this
effect, but the BM25 index itself still incorporates these terms.

\subsection{Static Corpus and Scraping Gap}

The current system operates over a frozen corpus extracted from PDF
archives of the 16th--18th Lok Sabhas.
The \texttt{QADownloader.py} and \texttt{QAScraper.py} scripts provide
infrastructure for downloading additional PDFs and metadata, but live
incremental ingestion into the running Qdrant collection is not implemented.
Newly enacted legislation or recent parliamentary sessions require a manual
re-ingestion run.

\subsection{Scanned Image Documents}

The \texttt{FailureAnalyzer.py} script catalogues a class of documents that
are physically present but contain no extractable digital text (scanned image
PDFs).
These documents---identified by extracted text length $<50$ characters---are
excluded from the corpus entirely.
Integrating an OCR layer (e.g., Tesseract or a cloud OCR API) would recover
this content but introduces additional processing cost and language complexity
for Devanagari-script scans.

% =============================================================================
\section{Conclusion and Future Work}
\label{sec:conclusion}

\subsection{Conclusion}

This paper presented a two-stage hybrid RAG system engineered under consumer-hardware
constraints to address three compounding failure modes in the Lok Sabha parliamentary
corpus: LLM hallucination of statutory facts, loss of numerical precision from
sparse keyword models, and multilingual cross-encoder degradation on Devanagari-script documents.

The core engineering contributions are:
(i)~an end-to-end data pipeline for constructing a structured, machine-readable
parliamentary corpus from scratch, covering scraping, PDF extraction, boilerplate
removal, failure cataloguing, and structural chunking across the 16th--18th
Lok Sabha archives;
(ii)~a structural regex chunker with a 400-word ceiling and parent-context
propagation that eliminates the blind-split loss of document context for both
Q\&A and Debate/Bill document types;
(iii)~a multilingual ingestion and retrieval strategy combining BAAI/bge-m3
(100+ languages) with the BAAI/bge-reranker-v2-m3 cross-encoder, mitigating
tokenisation-induced truncation on Devanagari-script passages; and
(iv)~a hardware configuration recipe enabling stable ingestion and querying of
a $\approx$200,000-document Qdrant vector database within a 16\,GB RAM envelope
via on-disk vector and HNSW storage.

The system establishes a replicable blueprint for deploying trustworthy,
precision-focused legal AI interfaces over institutionally important corpora
under real-world infrastructure constraints, building on validated design
principles from KemenkeuGPT~\cite{kemenkeu2024} and multilingual RAG
research~\cite{mrag_settings2024,multilingualrag_2025}.

\subsection{Future Work}

\textbf{GraphRAG for Legislative Amendment Chains}:
The amendment history of a single legislative Act spans multiple Lok Sabha
sessions; this temporal relationship cannot be recovered by chunk-level
retrieval alone.
Graph-based RAG approaches~\cite{graphrag2025,tkgrag2025,kgrag2025} would
construct a legislative knowledge graph, enabling traversal queries such as
``Show all amendments to Section~X of Act~Y across sessions.''

\textbf{OCR Integration for Scanned PDFs}:
Adding a Tesseract or cloud OCR layer to the \texttt{QAExtractor.py} pipeline
would recover the currently excluded scanned-image documents.
Devanagari OCR is particularly relevant for older Lok Sabha session records
where digital text is unavailable.

\textbf{RAGAS Quantitative Evaluation}:
Constructing a human-annotated benchmark of statutory fact queries with
gold-standard answers would enable systematic measurement of RAG faithfulness,
answer correctness, context precision, and context recall using the RAGAS
framework, analogous to the evaluation methodology of KemenkeuGPT~\cite{kemenkeu2024}.

\textbf{Corrective RAG (CRAG) with Agentic Evaluators}:
A CRAG loop could augment the current pipeline with an agent that scores
the Stage-2 cross-encoder output and, when the top-$k$ average score falls
below a confidence threshold, triggers a query expansion or broader index
scan before generation.

\textbf{Incremental Live Ingestion}:
Extending the pipeline to ingest a parliamentary RSS or Sansad TV metadata
feed would require streaming delta upload logic, incremental HNSW segment
merging under the Low-RAM Straitjacket constraints, and automated re-runs
of the language detection and chunking stages on newly downloaded PDFs.

% =============================================================================
\section*{Acknowledgment}
The authors gratefully acknowledge the Lok Sabha Secretariat for maintaining
publicly accessible archives of parliamentary proceedings and the developers
of Qdrant, BAAI/bge-m3, and the HuggingFace ecosystems.

% ─── References ──────────────────────────────────────────────────────────────
\begin{thebibliography}{00}

\bibitem{lewis2020rag}
  P.~Lewis \textit{et al.},
  ``Retrieval-Augmented Generation for Knowledge-Intensive NLP Tasks,''
  \textit{Advances in Neural Information Processing Systems (NeurIPS)},
  vol.~33, pp.~9459--9474, 2020.

\bibitem{rag_framework2025}
  Author(s) TBD,
  ``A Retrieval-augmented Generation Framework with Retriever and Generator
  Modules for Enhancing Factual Consistency,''
  \textit{ResearchGate}, 2025. [Online]. Available:
  \url{https://www.researchgate.net/publication/393590507}

\bibitem{graphrag2025}
  Author(s) TBD,
  ``To Enhance Graph-Based Retrieval-Augmented Generation (RAG) with Robust
  Retrieval Techniques,''
  in \textit{Proc. IEEE}, 2025, doi:~10.1109/ACCESS.2025.10871140.

\bibitem{tkgrag2025}
  Author(s) TBD,
  ``TKG-RAG: A Retrieval-Augmented Generation Framework with Text-chunk
  Knowledge Graph,''
  in \textit{Proc. IEEE}, 2025, doi:~10.1109/ACCESS.2025.10877117.

\bibitem{kgrag2025}
  Author(s) TBD,
  ``Knowledge Graph-Guided Retrieval Augmented Generation,''
  in \textit{Proc. 2025 Annual Conference of the North American Chapter of
  the ACL (NAACL)}, 2025, pp. [INSERT]. [Online]. Available:
  \url{https://aclanthology.org/2025.naacl-long.449}

\bibitem{kemenkeu2024}
  G.~F.~Febrian and G.~Figueredo,
  ``KemenkeuGPT: Leveraging a Large Language Model on Indonesia's Government
  Financial Data and Regulations to Enhance Decision Making,''
  \textit{arXiv preprint arXiv:2407.21459}, 2024.

\bibitem{multilingualrag_2025}
  L.~Ranaldi, B.~Haddow, and A.~Birch,
  ``Multilingual Retrieval-Augmented Generation for Knowledge-Intensive
  Task,''
  in \textit{Findings of EACL 2026}, arXiv:2504.03616, 2025.

\bibitem{mrag_settings2024}
  N.~Chirkova and V.~Nikoulina,
  ``Retrieval-augmented generation in multilingual settings,''
  \textit{arXiv preprint arXiv:2407.01463}, 2024. [Online]. Available:
  \url{https://arxiv.org/abs/2407.01463}

\end{thebibliography}

\balance
\end{document}
